\documentclass[12pt]{article}
\usepackage[utf8]{inputenc}
\usepackage[portuguese]{babel}
\usepackage{graphicx}
\usepackage{float}
\usepackage{geometry}
\usepackage{caption}
\usepackage{amsmath}
\captionsetup{labelfont=bf}
\geometry{a4paper, margin=2.5cm}

\title{Identificação de Dígitos através de Características Extraídas de Sinais de Áudio}
\author{Simão Tomás Botas Carvalho \\ nº2021223055}
\date{}

\begin{document}

\begin{figure}[H]
    \centering
    \includegraphics[width=0.8\textwidth]{uc_capa.png}
\end{figure}

\begingroup
\let\newpage\relax
\maketitle
\endgroup

\section{Introdução}

Hoje em dia, sistemas de reconhecimento de voz estão presentes em quase todos os dispositivos - de telemóveis a assistentes virtuais. Por trás desses sistemas está um processo essencial: a análise e extração de características dos sinais de áudio.

Neste projeto, o objetivo é identificar dígitos falados (de 0 a 9, em inglês) a partir de gravações de voz. Para isso, vamos trabalhar com sinais de áudio reais, explorando diferentes formas de os analisar — tanto no tempo como na frequência.

Ao longo do trabalho, vamos limpar e preparar os sinais, extrair características que ajudam a distinguir os dígitos, e por fim, usar essas informações para os reconhecer automaticamente. O dataset usado é o AudioMNIST, que inclui milhares de gravações feitas por diferentes pessoas.

Este relatório segue a estrutura das metas definidas no enunciado, refletindo passo a passo o desenvolvimento do projeto.

\section{Metodologia Geral}

Para o desenvolvimento deste trabalho foi usada a linguagem Matlab para tratamento de dados, assim como a criação de gráficos e plots para melhor demonstrar o processamento e a escolha de \textit{features} para a deteção de dígitos.

Foi desenvolvida uma estrutura de dados (tabela) que foi sendo atualizada ao longo do trabalho, como pedido no enunciado. Entre etapas, a estrutura é guardada em ficheiros \texttt{.mat}, permitindo reutilização e reduzindo tempo de execução.

\section{Meta 1 - Análise Temporal e de Frequência}

\subsection{Organização dos Dados}

O script principal começa por carregar todos os ficheiros \texttt{.wav} e guardar os metadados (diretório, participante, dígito, repetição, taxa de amostragem, etc.) juntamente com o áudio numa estrutura organizada (tabela).

\subsection{Importação e Visualização}

Um exemplo dos sinais originais é representado graficamente, com o eixo do tempo em segundos e as etiquetas correspondentes aos dígitos. Esta etapa permite observar visualmente os padrões dos sinais antes de qualquer modificação.

\begin{figure}[H]
    \centering
    \includegraphics[width=0.8\textwidth]{fig1.png}
    \caption{Exemplo de sinais de áudio originais para diferentes dígitos}
    \label{fig:sinais_originais}
\end{figure}

A geração de gráficos é opcional e pode ser ativada/desativada através de parâmetros do script, para melhorar o fluxo de execução quando necessário.

\subsection{Pré-processamento}

De seguida, é realizado o pré-processamento dos ficheiros de áudio:

\begin{itemize}
    \item Remoção do silêncio no início dos sinais, com base na energia calculada em pequenas janelas temporais e um limiar;
    \item Normalização da amplitude do áudio para o intervalo [-1,1];
    \item Equalização da duração de todos os sinais para um comprimento fixo de $L$ segundos (por defeito $L=0.5$ s), recorrendo a \textit{zero padding} quando o sinal é mais curto e truncagem quando é mais longo.

    \begin{figure}[H]
    \centering
    \includegraphics[width=0.8\textwidth]{fig2.png}
    \caption{Sinais após pré-processamento (remoção de silêncio, normalização e padding)}
    \label{fig:sinais_processados}
    \end{figure}
\end{itemize}

\subsection{Extração de Características Temporais}

Nesta secção são extraídas várias características temporais, com um particionamento do sinal em $D$ divisões (por defeito $D=3$).

Considere-se um sinal pré-processado $x[n]$, com $n=0,\dots,N-1$. O sinal é dividido em $D$ segmentos contíguos. Para a divisão $d\in\{1,\dots,D\}$, define-se o conjunto de índices $S_d$ correspondente a essa janela.

\begin{itemize}
    \item \textbf{Energia total}:
    \[
    E_{\text{tot}} = \sum_{n=0}^{N-1} x[n]^2
    \]

    \item \textbf{Energia por divisão}:
    \[
    E_d = \sum_{n\in S_d} x[n]^2,\qquad d=1,\dots,D
    \]

    \item \textbf{Amplitude máxima por divisão}:
    \[
    A_d = \max_{n\in S_d} |x[n]|,\qquad d=1,\dots,D
    \]

    \item \textbf{Resumo por amostra (média das divisões)} (quando aplicável nas visualizações):
    \[
    \bar{E} = \frac{1}{D}\sum_{d=1}^{D} E_d,
    \qquad
    \bar{A} = \frac{1}{D}\sum_{d=1}^{D} A_d
    \]

    \item \textbf{Desvio padrão (por dígito)}: para cada dígito, calcula-se o desvio padrão da energia total e da amplitude máxima (por divisão), considerando todas as amostras desse dígito:
\[
\sigma_{E} = \sqrt{\frac{1}{M}\sum_{i=1}^{M}(E_{\text{tot}}^{(i)}-\mu_{E})^2},
\qquad
\sigma_{A_d} = \sqrt{\frac{1}{M}\sum_{i=1}^{M}(A_d^{(i)}-\mu_{A_d})^2}
\]
\end{itemize}

Para a comparação e visualização inicial (incluindo um gráfico 3D), foram privilegiadas três características simples e informativas:

\begin{itemize}
    \item Energia total do sinal;
        \begin{figure}[H]
            \centering
            \includegraphics[width=0.8\textwidth]{energia_total.png}
            \caption{Energia total do sinal por dígito}
            \label{fig:energia_total}
        \end{figure}

    \item Amplitude máxima por divisão
        \begin{figure}[H]
            \centering
            \includegraphics[width=0.8\textwidth]{amplitude_max.png}
            \caption{Amplitude máxima por divisão e por dígito}
            \label{fig:amplitude_max_div}
        \end{figure}

    \item Energia por divisão
        \begin{figure}[H]
            \centering
            \includegraphics[width=0.8\textwidth]{energia_divisao.png}
            \caption{Energia por divisão}
            \label{fig:energia_por_div}
        \end{figure}
\end{itemize}

Para comparação das 3 métricas, criei um plot 3d para melhor visualização:
    \begin{figure}[H]
            \centering
            \includegraphics[width=1\textwidth]{plot3d.png}
            \caption{Comparação das 3 Métricas Temporais}
            \label{fig:plot3d}
    \end{figure}


\subsubsection{Observações e escolha qualitativa de \textit{features}}

Com base nos gráficos obtidos (por exemplo, Figuras \ref{fig:energia_total}, \ref{fig:amplitude_max_div} e \ref{fig:energia_por_div}), é possível perceber quais as características mais úteis.

A \textbf{energia total} representa a intensidade global do sinal. Ajuda a distinguir alguns dígitos, mas por si só não é suficiente, já que existe bastante sobreposição entre classes.

A \textbf{energia por divisão} já permite perceber melhor como a energia está distribuída ao longo do tempo (início, meio e fim do sinal), o que ajuda a diferenciar dígitos com comportamentos semelhantes.

A \textbf{amplitude máxima por divisão} destaca picos locais do sinal, sendo útil para capturar variações rápidas e complementar a informação da energia.

Estas características foram então guardadas na estrutura e num ficheiro \texttt{.mat}. Foi também criada uma versão mais leve da estrutura, sem os sinais (e, quando aplicável, sem coeficientes de Fourier), para acelerar o carregamento nas metas seguintes.

\subsection{Transformada de Fourier}

Para cada sinal de áudio é aplicada a Transformada de Fourier:

\[
X(f) = \sum_{n=0}^{N-1} x(n)e^{-j2\pi fn/N}
\]

Esta transformação permite analisar o sinal no domínio da frequência, mostrando como a energia está distribuída pelas diferentes frequências. Para a análise, considera-se apenas a parte positiva do espectro (até à frequência de Nyquist) e o espectro de amplitude normalizado.

\subsection{Extração de Estatísticas Espectrais}

A partir da Transformada de Fourier, são calculadas estatísticas espectrais agregadas por dígito, considerando todas as amostras disponíveis de cada classe:

\begin{itemize}
    \item Espectro de amplitude mediano;
    \item Primeiro quartil (Q1, 25\%);
    \item Terceiro quartil (Q3, 75\%).
\end{itemize}

Estas estatísticas são usadas para visualizar diferenças típicas no conteúdo espectral entre dígitos.

\begin{figure}[H]
            \centering
            \includegraphics[width=1\textwidth]{stats_spectral.png}
            \caption{Estatísticas Espetrais}
            \label{fig:stats_spec}
        \end{figure}

\subsection{Características Espectrais Calculadas}
Para cada amostra, foram calculadas várias características espectrais com base no espectro de amplitude $|X[k]|$.

Para isso, considera-se também a potência espectral:
\[
P[k] = |X[k]|^2,
\]
utilizando apenas as frequências não-negativas.

As principais características calculadas foram:

\begin{itemize}
    \item \textbf{Pico de frequência} (\textit{PicoFreq}) e \textbf{amplitude no pico} (\textit{PicoAmp}):
    \[
    k_{\max}=\arg\max_k |X[k]|,\qquad \textit{PicoFreq}=f_{k_{\max}},\qquad \textit{PicoAmp}=|X[k_{\max}]|.
    \]

    \item \textbf{Amplitude média do espectro} (\textit{MediaAmp}):
    \[
    \textit{MediaAmp} = \frac{1}{K}\sum_{k=1}^{K} |X[k]|
    \]
    onde $K$ é o número de bins considerados (frequências não-negativas).

    \item \textbf{\textit{Spectral Edge Frequency} a 95\%} (\textit{SpectralEdge}): define-se a potência acumulada
    \[
    C[m]=\sum_{k=1}^{m} P[k],\qquad C_{\text{tot}}=\sum_{k=1}^{K}P[k].
    \]
    Então, \textit{SpectralEdge} é a menor frequência $f_m$ tal que
    \[
    \frac{C[m]}{C_{\text{tot}}}\ge 0.95.
    \]

    \item \textbf{Centroide espectral} (\textit{SpectralCentroid}):
    \[
    \textit{SpectralCentroid}=\frac{\sum_{k=1}^{K} f_k\,P[k]}{\sum_{k=1}^{K} P[k]}.
    \]

    \item \textbf{Largura de banda espectral} (\textit{SpectralBW}):
    \[
    \textit{SpectralBW}=\sqrt{\frac{\sum_{k=1}^{K}(f_k-\textit{SpectralCentroid})^2P[k]}{\sum_{k=1}^{K}P[k]}}.
    \]

    \item \textbf{Entropia espectral} (\textit{SpectralEntropy}): com a distribuição normalizada
    \[
    p[k]=\frac{P[k]}{\sum_{m=1}^{K}P[m]},
    \]
    define-se
    \[
    \textit{SpectralEntropy}=-\sum_{k=1}^{K} p[k]\log(p[k]).
    \]
\end{itemize}

Numa fase inicial de seleção qualitativa (por inspeção dos gráficos), destacaram-se como particularmente úteis para separação entre classes: o pico de frequência, a amplitude média e a \textit{spectral edge} a 95\%. As restantes características mantêm-se igualmente disponíveis na estrutura para comparação e utilização em etapas seguintes.

A comparação pode ser visualizada de melhor forma na imagem a baixo:

\begin{figure}[H]
            \centering
            \includegraphics[width=1\textwidth]{spectral_car.png}
            \caption{Características Espectrais}
            \label{fig:spectral_car}
        \end{figure}

\section{Conclusões}

Neste trabalho foi desenvolvido um \textit{pipeline} para análise de sinais de áudio com o objetivo de identificar dígitos falados. O processo incluiu organização dos dados, pré-processamento, extração de características temporais e análise no domínio da frequência com a Transformada de Fourier.

O \textbf{pré-processamento} foi essencial para reduzir variações irrelevantes, como silêncio inicial, diferenças de amplitude e duração dos sinais. Isto permitiu uma comparação mais consistente entre amostras.

Relativamente às \textbf{características temporais}, a energia total revelou-se pouco discriminativa por si só, devido à sobreposição entre dígitos. Já a \textbf{energia por divisão} e a \textbf{amplitude máxima por divisão} mostraram melhores resultados, por captarem a evolução do sinal ao longo do tempo.

No domínio da frequência, as \textbf{características espectrais} permitiram observar padrões mais estáveis. Destacaram-se o \textbf{pico de frequência}, a \textbf{amplitude média} e a \textbf{\textit{spectral edge} a 95\%}, que ajudam a diferenciar os dígitos.

No geral, a combinação de características no tempo e na frequência dá uma representação mais completa do sinal. Na meta seguinte, estas \textit{features} serão usadas no classificador, permitindo avaliar o desempenho do sistema através de métricas como \textit{accuracy} e matriz de confusão.

\end{document}